\documentclass[10pt]{extarticle}
\usepackage[a4paper,margin=1.5cm]{geometry}
\usepackage[T1]{fontenc}
\usepackage{amsmath,amssymb}
\usepackage{graphicx}
\usepackage[export]{adjustbox}
\usepackage{multicol}
\usepackage{enumitem}
\usepackage{booktabs,array,tabularx}
\setlength{\columnsep}{0.7cm}
\setlength{\parindent}{0pt}
\setlength{\parskip}{3pt}
\pagestyle{empty}
\makeatletter
% figure registry (figdefs.tex is generated by the extract step)
\newcommand{\qfigdef}[4]{\@namedef{qf@file@#1}{#2}\@namedef{qf@w@#1}{#3}\@namedef{qf@wide@#1}{#4}}
\newcommand{\qf@check}[1]{\@ifundefined{qf@file@#1}{\PackageError{score2}{Unknown figure id #1}{}}{}}
\newcommand{\qf@img}[4][]{\includegraphics[width=#2\dimexpr\@nameuse{qf@w@#4}pt\relax,max width=#3,max height=0.45\textheight,#1]{\@nameuse{qf@file@#4}}}
% Full-width questions: finalize wraps every question that contains a wide figure in \qfullbegin...\qfullend,
% so the WHOLE question (stem, figures, options) is set across the page in reading order, then the two columns
% resume. The question's head - its number, stem and first figure (\qfullsplit marks the end of the head) -
% is kept on one page, so a number is never left at a page foot with its figure overleaf; the rest may break
% across pages normally, so no big gaps are left. Without a \qfullsplit, the whole question is kept together
% if it is under 45% of a page. Never write these yourself.
\newif\ifqf@full
\newif\ifqf@boxed
\newsavebox{\qf@box}
\newcommand{\qf@open}{\global\qf@fulltrue\par\global\qf@boxedtrue
  \setbox\qf@box\vbox\bgroup\hsize\textwidth\linewidth\textwidth}
\newcommand{\qf@emit}[1]{\par\egroup\global\qf@boxedfalse
  \ifdim\dimexpr\ht\qf@box+\dp\qf@box\relax<#1\textheight\noindent\usebox{\qf@box}\par\else\unvbox\qf@box\fi}
\newcommand{\qfullsplit}{\ifqf@boxed\qf@emit{0.85}\fi}
\newcommand{\qf@close}{\ifqf@boxed\qf@emit{0.45}\fi\global\qf@fullfalse}
\newcommand{\qfullbegin}{\end{multicols}\qf@open}
\newcommand{\qfullend}{\qf@close\begin{multicols}{2}}
\newcommand{\qfullstart}{\qf@open}
\newcommand{\qfulldone}{\qf@close}
% \qfig[scale]{id}: block figure, centred, directly where it is written (size set by the tool).
% Fallback if a wide figure is not inside a full-width question: break out of the columns in place - no floats.
\newcommand{\qfig}[2][1]{\qf@check{#2}%
  \ifqf@full\qf@blk{#1}{#2}\else
  \ifnum\@nameuse{qf@wide@#2}=1 \expandafter\@firstoftwo\else\expandafter\@secondoftwo\fi
  {\end{multicols}\begin{center}\qf@img{#1}{\textwidth}{#2}\end{center}\begin{multicols}{2}}%
  {\qf@blk{#1}{#2}}\fi}
\newcommand{\qf@blk}[2]{\par\smallskip{\centering\qf@img{#1}{\linewidth}{#2}\par}\smallskip}
% \qfigi[scale]{id}: inline figure (inside option lists / tables / minipages); never breaks columns.
\newcommand{\qfigi}[2][1]{\qf@check{#2}\qf@img[valign=c]{#1}{\linewidth}{#2}}
\makeatother
% structure macros - lint reads these to check completeness and figure placement
\newcommand{\papersec}[1]{\par\bigskip{\centering\Large\bfseries PAPER-#1\par}\medskip}
\newcommand{\partsec}[2]{\par\medskip\noindent{\large\bfseries PART-#1: #2}\par\smallskip}
\newcommand{\sect}[1]{\par\smallskip\noindent\textbf{#1}\par}
\newcommand{\qn}[1]{\par\smallskip\noindent\textbf{#1.}~}
\newcommand{\paraq}[2]{\par\smallskip\noindent\textbf{Paragraph for Questions #1 and #2}\par}
\newenvironment{qoptions}{\begin{itemize}[leftmargin=1.8em,itemsep=1pt,topsep=2pt,label={}]}{\end{itemize}}
\newcommand{\opt}[2]{\item[(#1)] #2}
% figure registry (generated by score2_tools.py finalize)
\qfigdef{p1_phy_s1_q02_1}{TOPS - Test 05 (14 Jan 2023)-img/p1_phy_s1_q02_1.png}{84.6}{0}
\qfigdef{p1_phy_s1_q03_1}{TOPS - Test 05 (14 Jan 2023)-img/p1_phy_s1_q03_1.png}{124.2}{0}
\qfigdef{p1_phy_s1_q04_1}{TOPS - Test 05 (14 Jan 2023)-img/p1_phy_s1_q04_1.png}{89.4}{0}
\qfigdef{p1_phy_s1_q06_1}{TOPS - Test 05 (14 Jan 2023)-img/p1_phy_s1_q06_1.png}{108.2}{0}
\qfigdef{p1_phy_s2_q07_1}{TOPS - Test 05 (14 Jan 2023)-img/p1_phy_s2_q07_1.png}{154.3}{0}
\qfigdef{p1_phy_s2_q08_1}{TOPS - Test 05 (14 Jan 2023)-img/p1_phy_s2_q08_1.png}{148.0}{0}
\qfigdef{p1_phy_s2_q09_1}{TOPS - Test 05 (14 Jan 2023)-img/p1_phy_s2_q09_1.png}{118.7}{0}
\qfigdef{p1_phy_s2_q10_1}{TOPS - Test 05 (14 Jan 2023)-img/p1_phy_s2_q10_1.png}{213.8}{0}
\qfigdef{p1_phy_s3_q02_1}{TOPS - Test 05 (14 Jan 2023)-img/p1_phy_s3_q02_1.png}{190.3}{0}
\qfigdef{p1_phy_s3_q03_1}{TOPS - Test 05 (14 Jan 2023)-img/p1_phy_s3_q03_1.png}{47.5}{0}
\qfigdef{p1_phy_s3_q04_1}{TOPS - Test 05 (14 Jan 2023)-img/p1_phy_s3_q04_1.png}{103.2}{0}
\qfigdef{p1_phy_s3_q07_1}{TOPS - Test 05 (14 Jan 2023)-img/p1_phy_s3_q07_1.png}{103.0}{0}
\qfigdef{p1_chem_s1_q03_A}{TOPS - Test 05 (14 Jan 2023)-img/p1_chem_s1_q03_A.png}{193.2}{0}
\qfigdef{p1_chem_s1_q03_B}{TOPS - Test 05 (14 Jan 2023)-img/p1_chem_s1_q03_B.png}{105.8}{0}
\qfigdef{p1_chem_s1_q03_3}{TOPS - Test 05 (14 Jan 2023)-img/p1_chem_s1_q03_3.png}{95.8}{0}
\qfigdef{p1_chem_s1_q03_C}{TOPS - Test 05 (14 Jan 2023)-img/p1_chem_s1_q03_C.png}{172.9}{0}
\qfigdef{p1_chem_s1_q03_D}{TOPS - Test 05 (14 Jan 2023)-img/p1_chem_s1_q03_D.png}{154.3}{0}
\qfigdef{p1_chem_s1_q04_A}{TOPS - Test 05 (14 Jan 2023)-img/p1_chem_s1_q04_A.png}{44.1}{0}
\qfigdef{p1_chem_s1_q06_A}{TOPS - Test 05 (14 Jan 2023)-img/p1_chem_s1_q06_A.png}{120.0}{0}
\qfigdef{p1_chem_s1_q06_2}{TOPS - Test 05 (14 Jan 2023)-img/p1_chem_s1_q06_2.png}{119.0}{0}
\qfigdef{p1_chem_s2_q09_1}{TOPS - Test 05 (14 Jan 2023)-img/p1_chem_s2_q09_1.png}{50.4}{0}
\qfigdef{p1_chem_s2_q09_2}{TOPS - Test 05 (14 Jan 2023)-img/p1_chem_s2_q09_2.png}{105.6}{0}
\qfigdef{p1_chem_s2_q09_3}{TOPS - Test 05 (14 Jan 2023)-img/p1_chem_s2_q09_3.png}{50.4}{0}
\qfigdef{p1_chem_s2_q09_4}{TOPS - Test 05 (14 Jan 2023)-img/p1_chem_s2_q09_4.png}{67.1}{0}
\qfigdef{p1_chem_s3_q03_1}{TOPS - Test 05 (14 Jan 2023)-img/p1_chem_s3_q03_1.png}{76.8}{0}
\qfigdef{p1_chem_s3_q04_1}{TOPS - Test 05 (14 Jan 2023)-img/p1_chem_s3_q04_1.png}{60.7}{0}
\qfigdef{p1_chem_s3_q04_2}{TOPS - Test 05 (14 Jan 2023)-img/p1_chem_s3_q04_2.png}{65.8}{0}
\qfigdef{p1_chem_s3_q04_3}{TOPS - Test 05 (14 Jan 2023)-img/p1_chem_s3_q04_3.png}{187.5}{0}
\qfigdef{p1_chem_s3_q05_1}{TOPS - Test 05 (14 Jan 2023)-img/p1_chem_s3_q05_1.png}{108.2}{0}
\qfigdef{p1_chem_s3_q06_1}{TOPS - Test 05 (14 Jan 2023)-img/p1_chem_s3_q06_1.png}{109.9}{0}
\qfigdef{p1_chem_s3_q06_2}{TOPS - Test 05 (14 Jan 2023)-img/p1_chem_s3_q06_2.png}{109.9}{0}
\qfigdef{p1_chem_s3_q06_3}{TOPS - Test 05 (14 Jan 2023)-img/p1_chem_s3_q06_3.png}{97.9}{0}
\qfigdef{p1_chem_s3_q06_4}{TOPS - Test 05 (14 Jan 2023)-img/p1_chem_s3_q06_4.png}{97.9}{0}
\qfigdef{p1_chem_s3_q06_5}{TOPS - Test 05 (14 Jan 2023)-img/p1_chem_s3_q06_5.png}{100.6}{0}
\qfigdef{p1_chem_s3_q06_6}{TOPS - Test 05 (14 Jan 2023)-img/p1_chem_s3_q06_6.png}{95.1}{0}
\qfigdef{p1_chem_s3_q06_7}{TOPS - Test 05 (14 Jan 2023)-img/p1_chem_s3_q06_7.png}{97.9}{0}
\qfigdef{p2_phy_s1_q01_1}{TOPS - Test 05 (14 Jan 2023)-img/p2_phy_s1_q01_1.png}{68.9}{0}
\qfigdef{p2_phy_s1_q02_1}{TOPS - Test 05 (14 Jan 2023)-img/p2_phy_s1_q02_1.png}{81.8}{0}
\qfigdef{p2_phy_s1_q03_1}{TOPS - Test 05 (14 Jan 2023)-img/p2_phy_s1_q03_1.png}{108.3}{0}
\qfigdef{p2_phy_s1_q04_1}{TOPS - Test 05 (14 Jan 2023)-img/p2_phy_s1_q04_1.png}{108.3}{0}
\qfigdef{p2_phy_s2_q05_1}{TOPS - Test 05 (14 Jan 2023)-img/p2_phy_s2_q05_1.png}{108.3}{0}
\qfigdef{p2_phy_s2_q06_1}{TOPS - Test 05 (14 Jan 2023)-img/p2_phy_s2_q06_1.png}{108.2}{0}
\qfigdef{p2_phy_s2_q07_1}{TOPS - Test 05 (14 Jan 2023)-img/p2_phy_s2_q07_1.png}{116.3}{0}
\qfigdef{p2_phy_s2_q08_1}{TOPS - Test 05 (14 Jan 2023)-img/p2_phy_s2_q08_1.png}{108.2}{0}
\qfigdef{p2_phy_s2_q10_1}{TOPS - Test 05 (14 Jan 2023)-img/p2_phy_s2_q10_1.png}{123.9}{0}
\qfigdef{p2_phy_s3_q01_1}{TOPS - Test 05 (14 Jan 2023)-img/p2_phy_s3_q01_1.png}{68.8}{0}
\qfigdef{p2_phy_s3_q02_1}{TOPS - Test 05 (14 Jan 2023)-img/p2_phy_s3_q02_1.png}{112.9}{0}
\qfigdef{p2_phy_s3_q04_1}{TOPS - Test 05 (14 Jan 2023)-img/p2_phy_s3_q04_1.png}{100.4}{0}
\qfigdef{p2_phy_s3_q05_1}{TOPS - Test 05 (14 Jan 2023)-img/p2_phy_s3_q05_1.png}{108.3}{0}
\qfigdef{p2_phy_s3_q06_1}{TOPS - Test 05 (14 Jan 2023)-img/p2_phy_s3_q06_1.png}{84.6}{0}
\qfigdef{p2_phy_s3_q07_1}{TOPS - Test 05 (14 Jan 2023)-img/p2_phy_s3_q07_1.png}{97.9}{0}
\qfigdef{p2_phy_s3_q08_1}{TOPS - Test 05 (14 Jan 2023)-img/p2_phy_s3_q08_1.png}{58.2}{0}
\qfigdef{p2_chem_s1_q02_A}{TOPS - Test 05 (14 Jan 2023)-img/p2_chem_s1_q02_A.png}{64.2}{0}
\qfigdef{p2_chem_s1_q02_B}{TOPS - Test 05 (14 Jan 2023)-img/p2_chem_s1_q02_B.png}{76.8}{0}
\qfigdef{p2_chem_s1_q02_C}{TOPS - Test 05 (14 Jan 2023)-img/p2_chem_s1_q02_C.png}{60.5}{0}
\qfigdef{p2_chem_s1_q02_D}{TOPS - Test 05 (14 Jan 2023)-img/p2_chem_s1_q02_D.png}{65.8}{0}
\qfigdef{p2_chem_s2_q07_A}{TOPS - Test 05 (14 Jan 2023)-img/p2_chem_s2_q07_A.png}{66.3}{0}
\qfigdef{p2_chem_s2_q07_B}{TOPS - Test 05 (14 Jan 2023)-img/p2_chem_s2_q07_B.png}{50.2}{0}
\qfigdef{p2_chem_s2_q07_C}{TOPS - Test 05 (14 Jan 2023)-img/p2_chem_s2_q07_C.png}{34.6}{0}
\qfigdef{p2_chem_s2_q07_D}{TOPS - Test 05 (14 Jan 2023)-img/p2_chem_s2_q07_D.png}{74.0}{0}
\qfigdef{p2_chem_s2_q08_A}{TOPS - Test 05 (14 Jan 2023)-img/p2_chem_s2_q08_A.png}{108.0}{0}
\qfigdef{p2_chem_s2_q08_B}{TOPS - Test 05 (14 Jan 2023)-img/p2_chem_s2_q08_B.png}{120.5}{0}
\qfigdef{p2_chem_s2_q08_C}{TOPS - Test 05 (14 Jan 2023)-img/p2_chem_s2_q08_C.png}{107.8}{0}
\qfigdef{p2_chem_s2_q08_D}{TOPS - Test 05 (14 Jan 2023)-img/p2_chem_s2_q08_D.png}{55.4}{0}
\qfigdef{p2_chem_s2_q08_5}{TOPS - Test 05 (14 Jan 2023)-img/p2_chem_s2_q08_5.png}{42.3}{0}
\qfigdef{p2_chem_s3_q02_1}{TOPS - Test 05 (14 Jan 2023)-img/p2_chem_s3_q02_1.png}{112.0}{0}
\qfigdef{p2_chem_s3_q02_2}{TOPS - Test 05 (14 Jan 2023)-img/p2_chem_s3_q02_2.png}{112.0}{0}
\qfigdef{p2_chem_s3_q04_1}{TOPS - Test 05 (14 Jan 2023)-img/p2_chem_s3_q04_1.png}{220.1}{0}
\qfigdef{p2_chem_s3_q04_2}{TOPS - Test 05 (14 Jan 2023)-img/p2_chem_s3_q04_2.png}{217.2}{0}
\qfigdef{p2_chem_s3_q04_3}{TOPS - Test 05 (14 Jan 2023)-img/p2_chem_s3_q04_3.png}{45.1}{0}
\qfigdef{p2_chem_s3_q05_1}{TOPS - Test 05 (14 Jan 2023)-img/p2_chem_s3_q05_1.png}{130.8}{0}
\raggedcolumns
\makeatletter\@beginparpenalty=10000\makeatother
\newcolumntype{L}[1]{>{\raggedright\arraybackslash}p{#1}}
\begin{document}
\begin{center}{\Large\bfseries TOPS --- Test 05 (14 Jan 2023)}\end{center}
\begin{multicols}{2}
\raggedcolumns
\papersec{1}
\partsec{1}{PHYSICS}
\sect{SECTION-I(i)}
\qn{1} Two particles, $P$ of mass $2m$ and $Q$ of mass $m$, are subjected to mutual force of attraction and no other force acts on them. At $t=0$, $P$ is at rest at point $O$ and $Q$ is moving away from $O$ with a speed $5u$. At a later instant $t=T$ (before any collision has taken place), $Q$ is moving towards $O$ with speed $u$. Then
\begin{qoptions}
\opt{A}{momentum of particle $P$ at $t=T$ is zero.}
\opt{B}{momentum of particle $P$ at $t=T$ is $6mu$.}
\opt{C}{work done by the force of attraction during $0\le t\le T$ is $12mu^2$.}
\opt{D}{work done by the force of attraction during $0\le t\le T$ is $-3mu^2$.}
\end{qoptions}
\qn{2} Two uniform disc $A$ and $B$ having masses 1 kg and 2 kg and having equal radii connected (at centres) by a massless spring of spring constant $k=1000$ N/m are kept on sufficient rough horizontal surface. Initially system is released from rest and spring is compressed by 10 cm, then
\qfig{p1_phy_s1_q02_1}
\begin{qoptions}
\opt{A}{friction force on $A$ and $B$ are always equal and opposite.}
\opt{B}{friction force on $A$ and $B$ are depends on masses}
\opt{C}{speed of disc $A$ is $\sqrt{\dfrac{20}3}$ m/s, when spring in its natural length.}
\opt{D}{frictional force on disc $A$ is zero, when spring in its natural length.}
\end{qoptions}
\qn{3} In the given figure, all surface are frictionless and strings pulleys are massless, then
\qfig{p1_phy_s1_q03_1}
\begin{qoptions}
\opt{A}{acceleration of block $B$ is $g/2$.}
\opt{B}{acceleration of block $B$ is zero.}
\opt{C}{acceleration of block wedge is zero.}
\opt{D}{acceleration of block wedge is $g/2$.}
\end{qoptions}
\qn{4} $n$ moles of an ideal gas undergo an isothermal process at temperature $T$. $P$-$V$ graph of the process is as shown in the figure. A point $A(V_1,P_1)$ is located on the $P$-$V$ curve. Tangent at point $A$, cuts the $V$-axis at point $D$. $AO$ is the line joining the point $A$ to the origin $O$ of $PV$ diagram. Then,
\qfig{p1_phy_s1_q04_1}
\begin{qoptions}
\opt{A}{coordinates of points $D$ is $\left(\dfrac{3V_1}2,0\right)$}
\opt{B}{coordinates of points $D$ is $(2V_1,0)$}
\opt{C}{area of the triangle $AOD$ is $nRT$}
\opt{D}{area of the triangle $AOD$ is $\dfrac34nRT$}
\end{qoptions}
\qn{5} Two moles of a monatomic ideal gas undergoes a thermodynamic process $\dfrac{V^3}{T^2}=$ constant, if the temperature is raised by 300 K then
\begin{qoptions}
\opt{A}{work done by the gas is 400 $R$}
\opt{B}{change in internal energy is 900 $R$.}
\opt{C}{molar heat capacity of the gas for the process is 13/6 $R$}
\opt{D}{molar heat capacity of the gas for the process is 3/2 $R$.}
\end{qoptions}
\qn{6} A uniform disc of mass $M$ and radius $R$ is lifted using a string as shown in the figure. Then, considering no slipping between string and disc :-
\qfig{p1_phy_s1_q06_1}
\begin{qoptions}
\opt{A}{its linear acceleration is $g$ upward}
\opt{B}{its linear acceleration is $g$ downward}
\opt{C}{its angular acceleration is $\dfrac{2g}R$}
\opt{D}{its rate of change of angular momentum is $MgR$}
\end{qoptions}
\sect{SECTION-I(ii)}
\qn{7} Figure shows a siphon. It is a long pipe which is used to drain water from the reservoir at higher level to a reservoir at lower level regarding with the siphon match the following columns :
\qfig{p1_phy_s2_q07_1}
{\footnotesize\setlength{\tabcolsep}{3pt}
\begin{tabular}{|l|p{3.6cm}|l|l|}
\hline
 & List--I & & List--II\\
\hline
(I) & Pressure is more than atmospheric pressure at & (P) & 1\\
(II) & Pressure is less than atmospheric pressure at & (Q) & 2\\
(III) & Pressure is highest of all the five points at & (R) & 3\\
(IV) & Pressure is least of all the five points. & (S) & 4\\
 & & (T) & 5\\
\hline
\end{tabular}}
\begin{qoptions}
\opt{A}{I$\to$Q;II$\to$P,R;III$\to$S;IV$\to$T}
\opt{B}{I$\to$P;II$\to$Q,R,S;III$\to$P;IV$\to$S}
\opt{C}{I$\to$R;II$\to$P;III$\to$Q;IV$\to$S}
\opt{D}{I$\to$P,Q,R;II$\to$Q,S;III$\to$S;IV$\to$T}
\end{qoptions}
\qn{8} Two light rods of length 1m each are hinged together as shown in figure. Rod $AB$ makes an angle $\theta$ with vertical while rod $BC$ makes an angle $\phi$ with horizontal. End $C$ of rod $BC$ remains in contact with horizontal rod. $AB$ is rotated with constant with constant angular velocity $\omega=1$ rad/s in clockwise direction. At the instant when $\theta=30^\circ$ and $\phi=30^\circ$. Match the variable in list--I with values in list--II.
\qfig{p1_phy_s2_q08_1}
{\footnotesize\setlength{\tabcolsep}{3pt}
\begin{tabular}{|l|p{3.2cm}|l|l|}
\hline
 & List--I & & List--II\\
\hline
(I) & Angular velocity of rod $BC$ in rad/s & (P) & $\dfrac{3\sqrt3+1}{3\sqrt3}$\\
(II) & Velocity of block $D$ in m/s & (Q) & $\dfrac{\sqrt3-1}{\sqrt6}$\\
(III) & Magnitude of angular acceleration of rod $BC$ in rad/s$^2$ & (R) & $\dfrac1{\sqrt3}$\\
(IV) & Acceleration of point $B$ in m/s$^2$ & (S) & 1\\
 & & (T) & 0\\
\hline
\end{tabular}}
\begin{qoptions}
\opt{A}{I$\to$R;II$\to$R;III$\to$P;IV$\to$S}
\opt{B}{I$\to$P;II$\to$R;III$\to$S;IV$\to$Q}
\opt{C}{I$\to$P;II$\to$Q;III$\to$R;IV$\to$S}
\opt{D}{I$\to$Q;II$\to$S;III$\to$T;IV$\to$R}
\end{qoptions}
\qn{9} One mole of a monatomic gas is taken through a cycle $ABCDA$ as shown in the $P$-$V$ diagram. List--II give the characteristics involved in the cycle. Match them with each of the processes given in list--I.
\qfig{p1_phy_s2_q09_1}
{\footnotesize\setlength{\tabcolsep}{3pt}
\begin{tabular}{|l|p{2.3cm}|l|p{2.3cm}|}
\hline
 & List--I & & List--II\\
\hline
(I) & Process $A\to B$ & (P) & Internal energy decreases\\
(II) & Process $B\to C$ & (Q) & Internal energy increases\\
(III) & Process $C\to D$ & (R) & Heat is lost\\
(IV) & Process $D\to A$ & (S) & Heat is gained\\
 & & (T) & Work is done on the gas\\
\hline
\end{tabular}}
\begin{qoptions}
\opt{A}{I$\to$P,R,T;II$\to$P,R;III$\to$Q,S;IV$\to$R,T}
\opt{B}{I$\to$P;II$\to$Q;III$\to$R,S,T;IV$\to$Q,T}
\opt{C}{I$\to$P,R,T;II$\to$P,R;III$\to$R,T;IV$\to$R,T}
\opt{D}{I$\to$R,T;II$\to$Q,S;III$\to$P;IV$\to$P,R,T}
\end{qoptions}
\qn{10} A copper rod (initially at room temperature of 20$^\circ$C) of non uniform cross section is placed between a steam chamber at 100$^\circ$C and Ice-water chamber at 0$^\circ$C.
\qfig{p1_phy_s2_q10_1}
{\footnotesize\setlength{\tabcolsep}{3pt}
\begin{tabular}{|l|p{3.2cm}|l|p{2cm}|}
\hline
 & List--I & & List--II\\
\hline
(I) & Initially rate of heat flow $\left(\dfrac{dQ}{dt}\right)$ will be & (P) & Maximum at section $A$\\
(II) & At steady state rate of heat flow $\left(\dfrac{dQ}{dt}\right)$ will be & (Q) & Maximum at section $B$.\\
(III) & At steady state temperature gradient $\left|\dfrac{dT}{dx}\right|$ will be & (R) & Maximum at section $C$.\\
(IV) & At steady state rate of change of temperature $\left(\dfrac{dT}{dt}\right)$ at a certain point will be. & (S) & Minimum at section $B$.\\
 & & (T) & Same for all sections.\\
\hline
\end{tabular}}
\begin{qoptions}
\opt{A}{I$\to$P;II$\to$T;III$\to$Q;IV$\to$T}
\opt{B}{I$\to$P;II$\to$S;III$\to$Q;IV$\to$T}
\opt{C}{I$\to$P;II$\to$T;III$\to$T;IV$\to$S}
\opt{D}{I$\to$P;II$\to$T;III$\to$Q;IV$\to$S}
\end{qoptions}
\sect{SECTION-II}
\qn{1} A car starts moving rectilinearly, first with acceleration $\alpha=5$ m/s$^2$ (the initial velocity is equal to zero), then uniformly, and finally, decelerating at the same rate $\alpha$ comes to a stop. The total time of motion equals $\tau=25$ s. The average velocity during that time is equal to $\langle v\rangle=72$ km/h. How much time (in sec) does the car move uniformly ?
\qn{2} Two identical blocks each having a mass of 20 kg are connected to each other by a light inextensible string as shown and are placed over a rough surface. Pulleys are connected to the blocks. Find acceleration $a_1$ (in m/s$^2$) of the blocks after one second after the application of the time varying force of $40\,t$ N, where $t$ is in second.
\qfig{p1_phy_s3_q02_1}
\qn{3} A light inextensible string that goes over a smooth fixed pulley as shown in the figure connects two blocks of masses 0.36 kg and 0.72 kg. Taking $g=10$ m/s$^2$, find the work done (in joules) by the string on block of mass 0.36 kg during the first second after the system is released from rest. ($g=10$ m/s$^2$)
\qfig{p1_phy_s3_q03_1}
\qn{4} A solid sphere has linear velocity $v_0=4$ m/s and angular velocity $\omega_0=9$ rad/s as shown. Ground on which it is moving is smooth. It collides elastically with rough wall of coefficient of friction $\mu$. Radius of sphere 1m and mass 2kg. Then what is coefficient of friction if sphere after collides rolls without slipping is.
\qfig{p1_phy_s3_q04_1}
\qn{5} How many grams of ice at $-14^\circ$C are needed to cool 200 g of water from $25^\circ$C to $10^\circ$C ? Take specific heat of ice $=0.5$ cal/g/$^\circ$C and latent heat of ice $=80$ cal/g.
\qn{6} One mole of a monatomic ideal gas undergoes an adiabatic expansion in which its volume becomes eight times its initial value. If the initial temperature of the gas is 100 K and the universal gas constant $R=8.0$ J mol$^{-1}$ K$^{-1}$, the decrease in its internal energy. in Joule, is \underline{\hspace{1.5cm}}
\qn{7} Two small beads each of mass $m$ are attached to a smooth ring of mass $m$ as shown in the figure. The whole system is kept on smooth horizontal table. Initially, each bead has given velocity $v_0$ as shown in figure. If the velocity of bead just before impact is $xv_0$, the value of $x$ is.
\qfig{p1_phy_s3_q07_1}
\qn{8} A railway track (made of iron) is laid in winter when the average temperature is 18$^\circ$C. The track consists of sections of 12.0 m placed one after the other. How much gap (in mm) should be left between two such sections so that there is no compression during summer when the maximum temperature goes to 48$^\circ$C ? Coefficient of linear expansion of iron $=11\times10^{-6}$ /$^\circ$C.
\partsec{2}{CHEMISTRY}
\sect{SECTION-I(i)}
\qn{1} The correct option(s) is/are :
\begin{qoptions}
\opt{A}{In 1 equivalent of $\mathrm{H_2SO_4}$, 2 mole of H-atom is present.}
\opt{B}{Volume strength of $\dfrac15$M - $\mathrm{H_2O_2}$ solution is `2.27 V'}
\opt{C}{14 gm of ethene required, 48 gm of oxygen gas for complete combustion.}
\opt{D}{Number of moles of oxygen atom in 276 gm of $\mathrm{N_2O_4}$ is 6.}
\end{qoptions}
\qn{2} For 1 mole of an ideal gas $\left(\dfrac{dP}{dT}\right)_V\cdot\left(\dfrac{dV}{dT}\right)_P=$
\begin{qoptions}
\opt{A}{$\dfrac RT$}
\opt{B}{$\dfrac1T$}
\opt{C}{$\dfrac{R^2}{PV}$}
\opt{D}{$\dfrac{R^2}{P^2}$}
\end{qoptions}
\qn{3} Identify compounds having $C_3$ axis of symmetry?
\begin{qoptions}
\opt{A}{\qfigi{p1_chem_s1_q03_A}}
\opt{B}{\qfigi{p1_chem_s1_q03_B}\par\qfigi{p1_chem_s1_q03_3}}
\opt{C}{\qfigi{p1_chem_s1_q03_C}}
\opt{D}{\qfigi{p1_chem_s1_q03_D}}
\end{qoptions}
\qn{4} Identify compounds having $C_3$ axis of symmetry?
\begin{qoptions}
\opt{A}{\qfigi{p1_chem_s1_q04_A}}
\opt{B}{Inorganic benzene}
\opt{C}{$\mathrm{PCl_5}$}
\opt{D}{Water}
\end{qoptions}
\qn{5} Two elements ``$A$'' and ``$X$'' of p-block form chlorides $\mathrm{ACl_3}$ and $\mathrm{XCl_3}$. The following information is given :

(i) Trifluoride of ``$A$'' is predominantly ionic and that of ``$X$'' is predominantly covalent.

(ii) Trihydride of ``$X$'' does not exist and that of ``$A$'' is electron deficient.

(iii) Both ``$A$'' and ``$X$'' exist as solid in nature with ``$X$'' existing as $\mathrm{X_2}$.

(iv) Period number of ``$A$'' is less than ``$X$'' in periodic table.

Which is/are correct :
\begin{qoptions}
\opt{A}{Dimeric $\mathrm{ACl_3}$ and $\mathrm{XCl_3}$ both are planar}
\opt{B}{$\mathrm{AX_3}$ is covalent in nature.}
\opt{C}{$\mathrm{ACl_3}$ is electron deficient but $\mathrm{A_2Cl_6}$ is electron precise molecule.}
\opt{D}{$\mathrm{XCl_3}$ is electron precise and $\mathrm{X_2Cl_6}$ is planar.}
\end{qoptions}
\qn{6} Which of the following is/are correct :
\begin{qoptions}
\opt{A}{\qfigi{p1_chem_s1_q06_A} has more value of partial pressure of $\mathrm{CO_2(P_{CO_2})}$ as compared to $\mathrm{P_{CO_2}}$ in \qfigi{p1_chem_s1_q06_2} at equilibrium (Assuming same temperature and volume)}
\opt{B}{Both $\mathrm{Ga_2Cl_6}$ and $\mathrm{Ga_2Cl_4}$ have Ga in +3 oxidation state only.}
\opt{C}{Benzoic acid dimerizes in Benzene.}
\opt{D}{On increasing polarization in ionic compounds, solubility generally increases in non-polar solvent.}
\end{qoptions}
\sect{SECTION-I(ii)}
\qn{7} {\footnotesize\setlength{\tabcolsep}{2pt}
\begin{tabular}{|l|p{3.1cm}|l|p{3.0cm}|}
\hline
 & List-I & & List-II\\
\hline
(P) & At Boyle's temperature in low pressure zone & (1) & $Z>1$\\
(Q) & At Boyle's temperature in high pressure zone & (2) & $Z=1$\\
(R) & Below Boyle's temperature in low pressure zone & (3) & Attractive forces are dominant between gas molecules\\
(S) & Above Boyle's temperature in high pressure zone & (4) & Repulsive forces are dominant between gas molecules\\
 & & (5) & $Z<1$\\
\hline
\end{tabular}}

Which is correctly matched :
\begin{qoptions}
\opt{A}{P$\to$5;Q$\to$1;R$\to$5;S$\to$1}
\opt{B}{P$\to$2;Q$\to$1;R$\to$5;S$\to$1}
\opt{C}{P$\to$3;Q$\to$1;R$\to$5;S$\to$2}
\opt{D}{P$\to$2;Q$\to$1;R$\to$4;S$\to$4}
\end{qoptions}
\qn{8} Which is incorrectly matched :

{\footnotesize\setlength{\tabcolsep}{2pt}
\begin{tabular}{|l|p{3.1cm}|l|p{3.0cm}|}
\hline
 & List-I & & List-II\\
\hline
(P) & At Boyle's temperature in low pressure zone & (1) & $Z>1$\\
(Q) & At Boyle's temperature in high pressure zone & (2) & $Z=1$\\
(R) & Below Boyle's temperature in low pressure zone & (3) & Attractive forces are dominant between gas molecules\\
(S) & Above Boyle's temperature in high pressure zone & (4) & Repulsive forces are dominant between gas molecules\\
 & & (5) & $Z<1$\\
\hline
\end{tabular}}

Which is incorrectly matched ?
\begin{qoptions}
\opt{A}{P$\to$2;Q$\to$1,4}
\opt{B}{R$\to$3,5;S$\to$1,4}
\opt{C}{P$\to$2;Q$\to$1,3}
\opt{D}{Q$\to$1,4;R$\to$3,5}
\end{qoptions}
\qn{9} {\footnotesize\setlength{\tabcolsep}{2pt}
\begin{tabular}{|l|p{3.1cm}|l|p{3.0cm}|}
\hline
 & List-I & & List-II\\
\hline
(P) & \qfigi{p1_chem_s2_q09_1} & (1) & Optically inactive compound\\
(Q) & \qfigi{p1_chem_s2_q09_2} & (2) & Optically active compound\\
(R) & \qfigi{p1_chem_s2_q09_3} & (3) & Dissymmetric compound\\
(S) & \qfigi{p1_chem_s2_q09_4} & (4) & Axis of symmetry present\\
 & & (5) & Asymmetric compound\\
\hline
\end{tabular}}
\begin{qoptions}
\opt{A}{P$\to$1;Q$\to$1,4;R$\to$2,5;S$\to$2,5}
\opt{B}{P$\to$2;Q$\to$3,4;R$\to$1,5;S$\to$2,5}
\opt{C}{P$\to$1,2;Q$\to$3,4;R$\to$4,5;S$\to$4}
\opt{D}{P$\to$3;Q$\to$2;R$\to$1;S$\to$1,2}
\end{qoptions}
\qn{10} Matching List :

{\footnotesize\setlength{\tabcolsep}{2pt}
\begin{tabular}{|l|p{3.1cm}|l|p{3.0cm}|}
\hline
 & List-I & & List-II\\
\hline
(P) & $\mathrm{N_xO_y}$ & (1) & For $x=2$; $y=5$, the specie is planar.\\
(Q) & $\mathrm{Cl_xO_y}$ & (2) & For $x=2$; $y=5$, the specie is non planar.\\
(R) & $\mathrm{(Si_xO_y)_n^{2n-}}$ if $n\in I$ & (3) & For $x=1$; $y=3$, specie contains $p\pi-p\pi$ bonds.\\
(S) & $\mathrm{S_xO_y}$ & (4) & For $x=1$; $y=2$, specie is polar.\\
 & & (5) & For $x=2$; $y=3$, specie is acidic in nature.\\
\hline
\end{tabular}}

Which of the following is a correct match?
\begin{qoptions}
\opt{A}{P$\to$1,3,5;Q$\to$3,4,5;R$\to$1,3;S$\to$3,4}
\opt{B}{P$\to$1,5;Q$\to$4,5;R$\to$2,3;S$\to$4,5}
\opt{C}{P$\to$1,4,5;Q$\to$4,5;R$\to$2;S$\to$3,4}
\opt{D}{P$\to$1,5;Q$\to$2,3,4;R$\to$2;S$\to$4}
\end{qoptions}
\sect{SECTION-II}
\qn{1} For 1 s orbital of hydrogen atom radial wave function is given as :

$R(r)=\dfrac1{\sqrt\pi}\left(\dfrac1{a_0}\right)^{3/2}e^{-r/a_0}$ (where $a_0=0.529$ \AA).

The ratio of radial probability density of finding electron at $r=a_0$ to the radial probability density of finding electron at the nucleus is given as $(x.e^{-y})$.

Calculate value of $(x+y)^3$.
\qn{2} In the chemical reaction between stoichiometric quantities of $\mathrm{KMnO_4}$ and $\mathrm{MnSO_4}$ in weakly alkaline solution, the number of moles of $\mathrm{MnO_2}$ formed is $Y$. If 4.2 moles of $\mathrm{KMnO_4}$ are consumed. The value of $Y$ is .......
\qn{3} Identify total number of stereo isomers for following structure
\qfig{p1_chem_s3_q03_1}
\qn{4} Identify total number of compounds which are more acidic than \qfigi{p1_chem_s3_q04_1}

\qfigi{p1_chem_s3_q04_2} (2) Me--COOH \quad (3) Me--SO$_3$H \quad (4) Picric acid

\qfig{p1_chem_s3_q04_3}
\qn{5} \qfig{p1_chem_s3_q05_1}
Total number of stereoisomers in the above compound is?
\qn{6} Number of molecular orbitals having ungerade symmetry among the following :

\qfig{p1_chem_s3_q06_1}
\qfig{p1_chem_s3_q06_2}
(3) \qfigi{p1_chem_s3_q06_3}

(4) \qfigi{p1_chem_s3_q06_4}

(5) \qfigi{p1_chem_s3_q06_5}

(6) \qfigi{p1_chem_s3_q06_6}

(7) \qfigi{p1_chem_s3_q06_7}
\qn{7} Number of species among the following having more ionic nature and less lattice energy compared to KF.

$\mathrm{LiF,\ KCl,\ CsF,\ BaI_2,\ NaF,\ NaBr,\ LiI,\ MgF_2}$
\qn{8} Number of species among the following which contain $sp^3$ hybridised atom involved in atleast 2 sigma bonds.

$\mathrm{BF_3}$, $\mathrm{N(SiH_3)_3}$, $\mathrm{O(CH_3)_2}$, $\mathrm{HN_3}$, $\mathrm{H_3BO_3}$, $\mathrm{H_3SiNCS}$, $\mathrm{PF_3}$, $\mathrm{N(CH_3)_3}$.
\partsec{3}{MATHEMATICS}
\sect{SECTION-I(i)}
\qn{1} If $a_{k+2}=a_k+\dfrac3{a_{k+1}}\ \forall\ k\in N$ where $a_1=1$, $a_2=3$, then which of the following is/are true ?
\begin{qoptions}
\opt{A}{$a_{100}\,a_{99}=300$}
\opt{B}{$a_{100}\,a_{99}=297$}
\opt{C}{$a_{2019}=\dfrac{2^{2018}}{{}^{2018}C_{1009}}$}
\opt{D}{$a_{2019}=\dfrac{2^{2019}}{{}^{2019}C_{1009}}$}
\end{qoptions}
\qn{2} Normals are drawn from the point $P(15,12)$ to the parabola $y^2=4x$. If the feet of the normals form a $\Delta ABC$, then which of the following is/are true
\begin{qoptions}
\opt{A}{Centroid of $\Delta ABC$ is $\left(\dfrac{26}3,0\right)$}
\opt{B}{Centroid of $\Delta ABC$ is $\left(\dfrac{22}3,0\right)$}
\opt{C}{Area of triangle formed by tangents to the parabola at $A$, $B$ and $C$ is 35 sq. units}
\opt{D}{Area of triangle formed by tangents to the parabola at $A$, $B$ and $C$ is 140 sq. units}
\end{qoptions}
\qn{3} If the polynomial $P(x)=24x^4+\lambda_1x^3+\lambda_2x^2+\lambda_3x+1$, where $\lambda_1,\lambda_2,\lambda_3\in R$ has four positive real roots $\alpha,\beta,\gamma,\delta$ such that $\alpha+2\beta+3\gamma+4\delta=4$, then
\begin{qoptions}
\opt{A}{$\lambda_1+\lambda_2+\lambda_3=-25$}
\opt{B}{$\lambda_1=-50$}
\opt{C}{$\lambda_3=-10$}
\opt{D}{$\lambda_2=25$}
\end{qoptions}
\qn{4} Let the conics $x^2=xy+1$ and $xy=x^2+1$ have eccentricities $e_1$ and $e_2$ respectively then
\begin{qoptions}
\opt{A}{$e_1>e_2$}
\opt{B}{$e_1<e_2$}
\opt{C}{$e_1\cdot e_2=2$}
\opt{D}{$e_1^{-2}+e_2^{-2}=1$}
\end{qoptions}
\qn{5} Let $n\in N$ ; $(2x^2-2x+1)^n=\displaystyle\sum_{r=0}^{2n}(-1)^ra_rx^r$ then
\begin{qoptions}
\opt{A}{$\displaystyle\sum_{r=0}^{2n}a_r=1$}
\opt{B}{$\displaystyle\sum_{r=1}^{2n}(-1)^ra_r=0$}
\opt{C}{$\displaystyle\sum_{k\geqslant0}{}^nC_k\,{}^{2n-2k}C_r=2^{n-r}a_r$ (wherever defined)}
\opt{D}{$\displaystyle\sum_{r=0}^{2n}\left((-1)^ra_ra_{2n-r}\right)=0$ if $n=$ odd}
\end{qoptions}
\qn{6} If a variable straight line passes through the point $P(2,0)$, intersects the curve $y^2=8x$ at the points $Q$, $R$ (let $PQ=\ell_1$ \& $PR=\ell_2$). Also given that $G$, $A$ represents geometric and arithmetic means of $\ell_1$ and $\ell_2$ respectively, then correct option(s) is/are
\begin{qoptions}
\opt{A}{$\ell_1,\ell_2\ge2$}
\opt{B}{$\ell_1,\ell_2<2$}
\opt{C}{$G\propto A$}
\opt{D}{$G\propto\sqrt A$}
\end{qoptions}
\sect{SECTION-I(ii)}
\qn{7} Match List-I with List-II and select the correct answer using the code given below the list.

{\footnotesize\setlength{\tabcolsep}{2pt}
\begin{tabular}{|l|p{3.9cm}|l|p{1.8cm}|}
\hline
 & List-I & & List-II\\
\hline
(I) & If the line $x-y-2=0$ is rotated through an angle $\dfrac\pi{12}$ in the clockwise direction about point $P(5,3)$ and equation of line in new position is $x-\sqrt3y+k=0$, then $k$ is & (P) & Rational\\
(II) & The point $R$ lying on $x$-axis such that $(PR+QR)$ is minimum, if $P$ and $Q$ are $(1,1)$ and $(3,2)$ respectively, then value of $(a+b)$ is [where $a$ and $b$ are coprime natural number] & (Q) & Irrational\\
(III) & Let $A(3,4)$, $B(0,0)$ and $C(3,0)$ be vertices of $\Delta ABC$. If $P(\alpha,\beta)$ is a point inside the $\Delta ABC$, equidistant from sides of triangle, then $(\alpha+\beta)$ is & (R) & Even integer\\
(IV) & Real value of $x$ satisfying $x^4-2x^3-x^2-2x+1=0$ is & (S) & Odd integer\\
 & & (T) & Prime number\\
\hline
\end{tabular}}
\begin{qoptions}
\opt{A}{I$\to$P;II$\to$P,R;III$\to$Q,S;IV$\to$P}
\opt{B}{I$\to$Q;II$\to$P,R;III$\to$P,S,T;IV$\to$Q}
\opt{C}{I$\to$Q;II$\to$Q,S;III$\to$Q;IV$\to$R}
\opt{D}{I$\to$P;II$\to$Q,S;III$\to$T;IV$\to$S}
\end{qoptions}
\qn{8} Match List-I with List-II and select the correct answer using the code given below the list.

{\footnotesize\setlength{\tabcolsep}{2pt}
\begin{tabular}{|l|p{4.4cm}|l|l|}
\hline
 & List-I & & List-II\\
\hline
(I) & If $\mathrm{cosec}\,\theta-\sin\theta=x^3$ and $\sec\theta-\cos\theta=y^3$ then $x^2y^2(x^2+y^2)$ is & (P) & 1\\
(II) & $\dfrac{\mathrm{cosec}\left(\frac{360^\circ}7\right)+\mathrm{cosec}\left(\frac{540^\circ}7\right)}{\mathrm{cosec}\left(\frac{180^\circ}7\right)}$ is & (Q) & 2\\
(III) & If $\cos(4x)=a_0+a_1\cos^2x+a_2\cos^4x$ is true for all values of $x\in R$, then $\dfrac{\left(a_0^2+a_1^2+a_2^2+a_0a_1a_2\right)}{13}$ is & (R) & 4\\
(IV) & $\left|\tan\dfrac\pi{16}+\tan\dfrac{5\pi}{16}+\right.$ $\left.\tan\dfrac{9\pi}{16}+\tan\dfrac{13\pi}{16}\right|$ is & (S) & 5\\
 & & (T) & 3\\
\hline
\end{tabular}}
\begin{qoptions}
\opt{A}{I$\to$P;II$\to$Q;III$\to$S;IV$\to$R}
\opt{B}{I$\to$P;II$\to$P;III$\to$S;IV$\to$R}
\opt{C}{I$\to$Q;II$\to$R;III$\to$S;IV$\to$P}
\opt{D}{I$\to$P;II$\to$Q;III$\to$R;IV$\to$S}
\end{qoptions}
\qn{9} Consider $P(x)=2x^3+x^2(3-7a)+x(3a^2-5a)+(3a^2+2a-1)$.

Match List-I with List-II and select the correct answer using the codes given below the list.

{\footnotesize\setlength{\tabcolsep}{2pt}
\begin{tabular}{|l|p{4.6cm}|l|l|}
\hline
 & List-I & & List-II\\
\hline
(I) & Number of positive value(s) of `$a$' for which $P(x)=0$ has exactly two equal roots, is & (P) & 3\\
(II) & Number of positive value(s) of `$a$' for which all roots of $P(x)=0$ are negative, is & (Q) & 1\\
(III) & Absolute value of sum of integral values of `$a$' for which one root of $P(x)=0$ is greater than $-1$ and other root is less than $-1$, is & (R) & 2\\
(IV) & Number of integral values of `$a$' for which all roots of $P(x)=0$ lie in $[-1,2]$, is & (S) & 0\\
 & & (T) & 4\\
\hline
\end{tabular}}
\begin{qoptions}
\opt{A}{I$\to$P;II$\to$S;III$\to$Q;IV$\to$R}
\opt{B}{I$\to$Q;II$\to$P;III$\to$S;IV$\to$R}
\opt{C}{I$\to$Q;II$\to$S;III$\to$P;IV$\to$R}
\opt{D}{I$\to$S;II$\to$P;III$\to$Q;IV$\to$R}
\end{qoptions}
\qn{10} $S\equiv x^2+y^2-2x-4y+1=0$ and $S'\equiv x^2+y^2-4x-2y+1=0$ are equations of two circles.

Match List-I with List-II and select the correct answer using the code given below the list.

{\footnotesize\setlength{\tabcolsep}{2pt}
\begin{tabular}{|l|p{4.6cm}|l|l|}
\hline
 & List-I & & List-II\\
\hline
(I) & Sum of number of tangents from $(1,3)$ to $S$ and number of tangents from $(1,3)$ to $S'$ is & (P) & 1\\
(II) & If $r$ is radius of circle passes through point of intersection of $S$ and $S'$ and whose normal is the straight line $2x+y-6=0$, then the value of $\dfrac{r^2}4$ is & (Q) & 2\\
(III) & If $m$ is slope of common tangent of $S=0$ and $S'=0$ then $|m|$ is & (R) & 3\\
(IV) & If $\ell$ is length of common tangent of $S$ and $S'$ then the value of $2\ell^2$ is & (S) & 4\\
 & & (T) & 7\\
\hline
\end{tabular}}
\begin{qoptions}
\opt{A}{I$\to$Q;II$\to$P;III$\to$Q;IV$\to$R}
\opt{B}{I$\to$Q;II$\to$Q;III$\to$P;IV$\to$S}
\opt{C}{I$\to$Q;II$\to$Q;III$\to$S;IV$\to$P}
\opt{D}{I$\to$P;II$\to$P;III$\to$S;IV$\to$Q}
\end{qoptions}
\sect{SECTION-II}
\qn{1} Given $a,b\in\{0,1,2,\ldots,6\}$. Consider the system of equations

$x+y+z=4$,

$2x+y+3z=6$,

$x+2y+az=b$

The number of ordered pairs $(a,b)$ so that the system of equations has unique solution, is
\qn{2} The number of 5-digit nutural numbers beginning with 9 that has exactly three identical digits (other digits are distinct) is
\qn{3} In the expansion of $(1-x)^{30}(1+x+x^2)^{31}$ the coefficient of $x^{47}$ is $\left(-{}^{30}C_r\right)$ then the value of $r$ is
\qn{4} In the equation $x^5-5x^4+9x^3-9x^2+5x-1=0$ if $p$ is the sum of real roots and $q$ is the product of real roots then the value of $\dfrac qp$ is
\qn{5} Let $\lambda$ be number of all possible 8 digit odd numbers formed by using only digits 0, 0, 2, 2, 3, 3, 4, 5, then $\dfrac\lambda{900}$ is
\qn{6} If $S$, $S_1$, $S_2$ be the circles of radii 5, 3, 2 respectively. If $S_1$ and $S_2$ touch externally and they touch internally with $S$. The radius of circle $S_3$ which touches externally with $S_1$ and $S_2$ and internally with $S$ is
\qn{7} The locus of point of intersection of the normals to the parabola $y^2=4ax$ drawn at the extremities of a chord subtending an angle of $90^\circ$ at origin is $y^2=ab(x-ac)$ then $\left(\dfrac bc\right)=$
\qn{8} A regular polygon of $n$ sides has 54 diagonals. The number of triangles formed with exactly one side of the triangle coinciding with the side of polygon, is
\papersec{2}
\partsec{1}{PHYSICS}
\sect{SECTION-I(i)}
\qn{1} A disc is rolling without sliding on a convex surface of same radius. Radius of curvature of topmost point $P$ at this instant is :-
\qfig{p2_phy_s1_q01_1}
\begin{qoptions}
\opt{A}{$3R$}
\opt{B}{$\dfrac{12R}5$}
\opt{C}{$\dfrac{8R}3$}
\opt{D}{None of these}
\end{qoptions}
\qn{2} In the figure shown, all surfaces are frictionless. A uniform rod is released from rest at $\theta=53^\circ$. The speed of the lower end at the moment $\theta=37^\circ$ is :-
\qfig{p2_phy_s1_q02_1}
\begin{qoptions}
\opt{A}{$\sqrt{\dfrac{g\ell}5}$}
\opt{B}{$\sqrt{\dfrac{72}{125}g\ell}$}
\opt{C}{$\sqrt{\dfrac{27}{125}g\ell}$}
\opt{D}{None of these}
\end{qoptions}
\qn{3} A pipe of radius $R$ and thickness $t$ ($\ll R$) is fixed at one end as shown in the figure. Length of the pipe is $\ell$ and modulus of rigidity $\eta$. Find torque required to produce a torsion of $\theta$.
\qfig{p2_phy_s1_q03_1}
\begin{qoptions}
\opt{A}{$\dfrac{2\pi R^3t\eta}{\ell}\theta$}
\opt{B}{$\dfrac{\pi R^3t\eta}{\ell}\theta$}
\opt{C}{$\dfrac{2\pi R^3t\eta}{3\ell}\theta$}
\opt{D}{$\dfrac{\pi R^3t\eta}{4\ell}\theta$}
\end{qoptions}
\qn{4} Figure shows a tank filled with water. Neglect viscosity and surface tension. Orifice area is very small.

$h_1$ = height of water column in tank

$h_2$ = height reached by fountain from orifice, then :-
\qfig{p2_phy_s1_q04_1}
\begin{qoptions}
\opt{A}{$h_1=h_2$}
\opt{B}{$h_1>h_2$}
\opt{C}{$h_1<h_2$}
\opt{D}{$\dfrac{h_1}{h_2}$ decreases with time}
\end{qoptions}
\sect{SECTION-I(ii)}
\qn{5} Choose the correct statement for Poisson's ratio of a body as shown in the figure.
\qfig{p2_phy_s2_q05_1}
\begin{qoptions}
\opt{A}{If $\sigma=0.5$, volume remain constant}
\opt{B}{If $\sigma<0.5$, volume increases}
\opt{C}{if $\sigma=0.2$ and axial strain is 2\%, percentage change in volume is 1.2\%.}
\opt{D}{None of these}
\end{qoptions}
\qn{6} An ideal gas is taken from thermodynamic state $A$ to $D$ as shown in the figure :-
\qfig{p2_phy_s2_q06_1}
\begin{qoptions}
\opt{A}{Total change in internal energy is zero.}
\opt{B}{Process $CD$ is isobaric.}
\opt{C}{Process $BC$ may be isothermal.}
\opt{D}{In process $CD$ heat is released.}
\end{qoptions}
\qn{7} The spectral emission power $E_\lambda$ for a body at temperature $T_1$ and temperature $T_2$ is plotted against the wavelength and area under the curve is found $A$ for $T_1$ and $9A$ for $T_2$.
\qfig{p2_phy_s2_q07_1}
\begin{qoptions}
\opt{A}{$\dfrac{T_1}{T_2}=\dfrac1{\sqrt3}$}
\opt{B}{$\dfrac{T_1}{T_2}=3$}
\opt{C}{$\dfrac{\lambda_1}{\lambda_2}=\sqrt3$}
\opt{D}{$\dfrac{\lambda_1}{\lambda_2}=\dfrac13$}
\end{qoptions}
\qn{8} A body of mass 2kg is moving on a straight line for which velocity-position ($v$-$x$) graph as shown in the figure. At $t=0$, body is at origin.
\qfig{p2_phy_s2_q08_1}
\begin{qoptions}
\opt{A}{Force acting on the body as $x=10$ m is 20 N.}
\opt{B}{$x=20\left(1-e^{-t}\right)$ m}
\opt{C}{Acceleration of the body is uniform.}
\opt{D}{Acceleration-position graph is parabola.}
\end{qoptions}
\qn{9} If force ($F$), momentum ($P$) and mass ($M$) are fundamental quantities in system of units, then dimensional formulae for length ($L$) and time ($T$) will be :-
\begin{qoptions}
\opt{A}{$[T]=\left[F^2P^{-1}M^0\right]$}
\opt{B}{$[L]=\left[F^{-1}M^{-2}P^{-1}\right]$}
\opt{C}{$[L]=\left[F^{-1}P^2M^{-1}\right]$}
\opt{D}{$[T]=\left[F^{-1}P^1M^0\right]$}
\end{qoptions}
\qn{10} A flat car is given an acceleration $a_0$ and started from rest as shown in the figure. Neglect friction between floor and car wheels and also the mass of the pulley. Mass of block $A=50$ kg. Then :-
\qfig{p2_phy_s2_q10_1}
\begin{qoptions}
\opt{A}{Tension in the string connected to block $A$ is 350 N.}
\opt{B}{Frictional force acting on $A$ is 150 N.}
\opt{C}{Acceleration of block $A$ is 4 m/s$^2$.}
\opt{D}{Acceleration of block $A$ is 1 m/s$^2$.}
\end{qoptions}
\sect{SECTION-III}
\qn{1} A uniform solid cylinder is released on a rough inclined surface of coefficient of friction $\mu=0.2$. If `$a$' is linear acceleration of centre of mass and $\alpha$ is angular acceleration, then find $\dfrac{11R\alpha}a$.
\qfig{p2_phy_s3_q01_1}
\qn{2} A container of mass 200 kg is moving towards right by applying a constant force 100 N. Initial velocity of container is zero. Sand is falling in container with constant rate of 1 kg/s with the help of a fixed cone as shown. Velocity of container at $t=200$ sec is $10N$ m/s. Find $N$. ($\ell$n 1.5 = 0.4)
\qfig{p2_phy_s3_q02_1}
\qn{3} A 1kg substance of specific heat capacity is given as $C=(A+Bt^2)$, where $A$ and $B$ are constant and $t$ is temperature in $^\circ$C, is heated. Magnitude of difference in average heat capacity $\left(\dfrac{\Delta Q}{\Delta t}\right)$ from $0^\circ$C to $t^\circ$C and heat capacity at $t^\circ$C is given as $\dfrac{4BT^2}{3n}$, find $n$.
\qn{4} Figure shows diagram of two bodies just before and just after collision. If $e$ is the coefficient of restitution, find $10\,e$.
\qfig{p2_phy_s3_q04_1}
\qn{5} Air is trapped in a horizontal glass tube by 36 cm Hg column as shown. If the tube is held vertically keeping the open end upside, length of the air column shrinks to 19 cm. The distance in cm by which mercury column shifts is :
\qfig{p2_phy_s3_q05_1}
\qn{6} A uniform disc of mass $M$ is free to rotate in vertical plane about a fixed horizontal axis passing through centre. The disc is at rest and a fly `$F$' of mass $m=M/2$ sits gently as the position shown in the diagram at $t=0$. The fly crawls on the disc in such a manner that its position with respect to ground remains fixed. At $t=2$ sec the insect stops crawling. If $v$ is the velocity of insect just after this instant. Find $v/5$. ($g=10$ m/s$^2$)
\qfig{p2_phy_s3_q06_1}
\qn{7} In the figure, the ring is placed in the mid point and released. If $M$ rises to a maximum height of 4m. Find $M/m$. Assume no friction between string and ring.
\qfig{p2_phy_s3_q07_1}
\qn{8} Figure shows a container of capacity 500 cc and cross-section area 20 cm$^2$, filled with 300 cc liquid at $20^\circ$C. Coefficient of volume expansion of liquid is $\gamma_L$ and coefficient of linear expansion of container is $\alpha_C$. Find $\dfrac{22\alpha_C}{\gamma_L}$ for which `$h$' remain constant at all temperature.
\qfig{p2_phy_s3_q08_1}
\partsec{2}{CHEMISTRY}
\sect{SECTION-I(i)}
\qn{1} A 1.5 m solution of acetic acid (I) is mixed with 3 m solution of the acetic acid (II) to prepare 2 m solution ($m$ is molality of solution). Select the correct statement(s) :

(I) Mass ratio of solvents mixed $\left(\dfrac{\mathrm I}{\mathrm{II}}\right)$ is $\dfrac12$

(II) Mass ratio of solvents mixed $\left(\dfrac{\mathrm I}{\mathrm{II}}\right)$ is $\dfrac21$

(III) Mass ratio of solutions mixed $\left(\dfrac{\mathrm I}{\mathrm{II}}\right)$ is $\dfrac{109}{59}$

(IV) Mass ratio of solutions mixed $\left(\dfrac{\mathrm I}{\mathrm{II}}\right)$ is $\dfrac{59}{109}$
\begin{qoptions}
\opt{A}{I, III}
\opt{B}{I, IV}
\opt{C}{II, III}
\opt{D}{II, IV}
\end{qoptions}
\qn{2} Which will not show tautomerism?
\begin{qoptions}
\opt{A}{\qfigi{p2_chem_s1_q02_A}}
\opt{B}{\qfigi{p2_chem_s1_q02_B}}
\opt{C}{\qfigi{p2_chem_s1_q02_C}}
\opt{D}{\qfigi{p2_chem_s1_q02_D}}
\end{qoptions}
\qn{3} Identify the specie having magnetic moment 0 B.M. and 6 electrons in bonding gerade molecular orbital :
\begin{qoptions}
\opt{A}{$\mathrm{Li_2}$}
\opt{B}{$\mathrm{B_2}$}
\opt{C}{$\mathrm{O_2^{-2}}$}
\opt{D}{$\mathrm{C_2}$}
\end{qoptions}
\qn{4} Choose the correct option :
\begin{qoptions}
\opt{A}{Boiling point : $\mathrm{SiCl_4>CCl_4}$}
\opt{B}{Number of $sp^3$ hybrid orbitals in Boric acid is 12.}
\opt{C}{Basicity of isohypophosphoric acid is 3.}
\opt{D}{Formula of 4-member single chain silicate and 4 member cyclic silicate is same.}
\end{qoptions}
\sect{SECTION-I(ii)}
\qn{5} A closed rigid container of volume 1 litre, initially contains 4 mole of HI(g) and 4.25 moles of $\mathrm{N_2(g)}$ at 27$^\circ$C. When the mixture is heated to 127$^\circ$C two reversible reactions takes place till equilibrium is established.

$\mathrm{2HI(g)\rightleftharpoons H_2(g)+I_2(g)}$ \quad $K_{C_1}$

$\mathrm{2N_2(g)+6H_2(g)\rightleftharpoons 4NH_3(g)}$ \quad $K_{C_2}$

At equilibrium the concentration of $\mathrm{I_2}$ and $\mathrm{NH_3}$ are 1.0 M and 0.5 M respectively. Choose the correct option :
\begin{qoptions}
\opt{A}{$K_{C_1}=\dfrac5{32}$}
\opt{B}{$K_{C_2}=\dfrac{16}{125}$}
\opt{C}{$K_{C_1}\times\sqrt{K_{C_2}}=\dfrac1{25}$}
\opt{D}{Conc. of $\mathrm{N_2}$ at equilibrium = 4.0 mol/litre}
\end{qoptions}
\qn{6} In the emission spectrum of H-atom from energy level `$n$' to ground state in one or more step, the wave length of lines belonging to Balmer series may be ($R$ = Rydberg constant) :
\begin{qoptions}
\opt{A}{$\dfrac{36}{5R}$}
\opt{B}{$\dfrac{16}{3R}$}
\opt{C}{$\dfrac1R$}
\opt{D}{$\dfrac{35}{3R}$}
\end{qoptions}
\qn{7} Which among the following carbocations is most stable?
\begin{qoptions}
\opt{A}{\qfigi{p2_chem_s2_q07_A}}
\opt{B}{\qfigi{p2_chem_s2_q07_B}}
\opt{C}{\qfigi{p2_chem_s2_q07_C}}
\opt{D}{\qfigi{p2_chem_s2_q07_D}}
\end{qoptions}
\qn{8} Among the following pair(s) of compounds identify enantiomeric pairs
\begin{qoptions}
\opt{A}{\qfigi{p2_chem_s2_q08_A}}
\opt{B}{\qfigi{p2_chem_s2_q08_B}}
\opt{C}{\qfigi{p2_chem_s2_q08_C}}
\opt{D}{\qfigi{p2_chem_s2_q08_D}\ \qfigi{p2_chem_s2_q08_5}}
\end{qoptions}
\qn{9} 24 gm of optically pure tartaric acid is dissolved in water to make 240 mL solution. It is kept in 20 cm polarimeter tube and plane polarised light is passed through it to produce rotation of $-2.4^\circ$. Calculate specific rotation for above compound?
\begin{qoptions}
\opt{A}{$12^\circ$}
\opt{B}{$-12^\circ$}
\opt{C}{$+6^\circ$}
\opt{D}{$-6^\circ$}
\end{qoptions}
\qn{10} Select the correct statement(s) about the compound $\mathrm{O_2[AsF_6]}$ :
\begin{qoptions}
\opt{A}{Magnetic moment of the species is zero}
\opt{B}{Oxidation number, hybridisation and geometry at As is +5, $sp^3d^2$ and octahedral respectively}
\opt{C}{It has $7\sigma$ and $1.5\pi$ bond}
\opt{D}{O--O bond order is same as in $\mathrm{KO_2}$}
\end{qoptions}
\sect{SECTION-III}
\qn{1} In absence of external electric and magnetic field and without considering the electronic spin, the degeneracy of the second excited state of H-atom is `$x$', degeneracy of the second excited state of $\mathrm{H^-}$ is `$y$'. And number of angular nodes in $d_{x^2-y^2}$ orbital is `$z$'. Calculate the value of $\dfrac xy+z$ ?
\qn{2} Equilibrium constant of $\mathrm{SO_3(g)}$ decomposition and formations was calculated by following two methods :
\qfig{p2_chem_s3_q02_1}
at $t=0$, 10 moles of $\mathrm{SO_3}$ is added 5($\ell$) in container. If at equilibrium, mixture is titrated with 0.2 M acidic $\mathrm{KMnO_4}$ its volume is 4 litre. $K_C$ for the reaction is `$x$' (mol/litre)
\qfig{p2_chem_s3_q02_2}
Initially 2 : 1 molar ratio $\mathrm{SO_2}$ and $\mathrm{O_2}$ is added respectively. At the equilibrium total pressure of mixture 20 atm. Mole \% of $\mathrm{O_2}$ at equilibrium is 25\%, $K_P$ for the reaction is `$Y$' atm$^{-1}$.

The value of $\dfrac yx$ is ?
\qn{3} A tiny particle moving with a velocity $\dfrac1{75}$th times the velocity of light. If uncertainty in the position is $\dfrac{6.626}{3.14}\times10^{-12}$m, then the minimum uncertainty in kinetic energy is $x\times10^{-16}$ J, then the value of $x=$ ?
\qn{4} How many compounds are soluble in aqueous NaOH?
\qfig{p2_chem_s3_q04_1}
\qfig{p2_chem_s3_q04_2}
(5) \qfigi{p2_chem_s3_q04_3}
\qn{5} Identify total number of chiral centers in following compound?
\qfig{p2_chem_s3_q05_1}
\qn{6} $X=$ Number of incorrect statement(s) among the following

$Y=$ Number of correct statement(s) among the following

Find $Y-X$ :

(i) Covalent character order : $\mathrm{ZnCl_2<CdCl_2<HgCl_2}$

(ii) Acidic nature of oxide : $\mathrm{CrO_3>CrO>Cr_2O_3}$

(iii) Solubility in water : $\mathrm{BeSO_4>MgSO_4>CaSO_4}$ $\mathrm{>SrSO_4>BaSO_4}$

(iv) Electrical conductivity in infinitely aqueous dilute solution : $\mathrm{CsCl>NaCl}$

(v) Thermal stability : $\mathrm{LiOH<NaOH<KOH}$ $\mathrm{<RbOH<CsOH}$

(vi) Intensity of colour : $\mathrm{AgCl>AgI}$

(vii) Solubility in acetone : $\mathrm{NaCl>NaI}$
\qn{7} `$X$' is the total number of atoms lying in a plane in $\mathrm{Al_2(CH_3)_6}$

`$Y$' is the number of species having unpaired electron present in hybrid orbital, among $\mathrm{NO_2}$, $\mathrm{ClO_2}$, $\mathrm{ClO_3}$, $\mathrm{CH_3}$, $\mathrm{CF_3}$.

Calculate : $X-2Y$.
\qn{8} Number of species among the following having less EA, more IE and more magnetic moment than O-atom.

B, Al, Si, F, Cl, $\mathrm{O^-}$, N, C, S, P.
\partsec{3}{MATHEMATICS}
\sect{SECTION-I(i)}
\qn{1} The value of the expression
\par\noindent\hfill$\dfrac{\sin146^\circ+\sin18^\circ+\sin196^\circ}{\sin3^\circ\sin57^\circ\sin63^\circ\sin73^\circ\sin98^\circ}$\hfill\null\par
is equal to
\begin{qoptions}
\opt{A}{16}
\opt{B}{4}
\opt{C}{12}
\opt{D}{8}
\end{qoptions}
\qn{2} If $(1-y)^{35}(1+y)^{45}=A_0+A_1y+A_2y^2+A_3y^3+\ldots..+A_{80}y^{80}$, then
\begin{qoptions}
\opt{A}{$\dfrac{A_2}{A_1}>2$}
\opt{B}{$A_1=A_2$}
\opt{C}{$\dfrac{A_2}{A_1}<1$}
\opt{D}{$1<\dfrac{A_2}{A_1}<2$}
\end{qoptions}
\qn{3} Let $A$ be the sum of the first 20 terms and $B$ be the sum of the first 40 terms of the series $1^3+2.2^3+3^3+2.4^3+5^3+2.6^3+\ldots..$ If $B-2A=600\lambda$, then $\lambda$ is equal to -
\begin{qoptions}
\opt{A}{1281}
\opt{B}{1381}
\opt{C}{1024}
\opt{D}{1481}
\end{qoptions}
\qn{4} The number of ways of arranging letters of the word COMPUTER such that none of the first three letters are $C$, $O$ or $M$, is
\begin{qoptions}
\opt{A}{7100}
\opt{B}{7000}
\opt{C}{7200}
\opt{D}{7120}
\end{qoptions}
\sect{SECTION-I(ii)}
\qn{5} If pair of tangents are drawn from centre $C$ of the circle $x^2+y^2-5x=0$ to the nearer branch of the hyperbola $\dfrac{x^2}{25}-\dfrac{y^2}{16}=1$ to touch it at point $A$ and $B$, then
\begin{qoptions}
\opt{A}{product of the slopes of lines $CA$ and $CB$ is $\dfrac{-64}{75}$}
\opt{B}{product of the slopes of lines $CA$ and $CB$ is $\dfrac{-32}{25}$}
\opt{C}{circumcentre of $\Delta ABC$ lies inside the $\Delta ABC$.}
\opt{D}{circumcentre of $\Delta ABC$ lies outside the $\Delta ABC$.}
\end{qoptions}
\qn{6} Let $A$ be the set containing value(s) of $x$ which satisfy the equation $3\sin^2x-\cos x=\sin x(3\cos x-1)$, where $x\in[0,9\pi]$, then mark the correct statement(s).
\begin{qoptions}
\opt{A}{Number of elements in set $A$ is 17}
\opt{B}{Number of elements in set $A$ is 18}
\opt{C}{Sum of elements in set $A$ is $\dfrac{297\pi}4$}
\opt{D}{Sum of elements in set $A$ is $\dfrac{295\pi}4$}
\end{qoptions}
\qn{7} A parabola touches straight lines $x-2y+13=0$ and $2x+y+1=0$ at points $P(17,15)$ and $Q(2,-5)$ respectively then :
\begin{qoptions}
\opt{A}{Focus of parabola is $\left(\dfrac72,-3\right)$}
\opt{B}{Directrix of parabola is $x+y-2=0$}
\opt{C}{Focus of parabola is $(5,-1)$}
\opt{D}{Length of latus rectum of parabola is 16}
\end{qoptions}
\qn{8} The triangle $ABC$ has $\angle A=60^\circ$, $\angle B=45^\circ$. The bisector of $\angle A$ meets the side $BC$ at $T$ where $AT=24$. Then which of the following is correct?
\begin{qoptions}
\opt{A}{length of $BC=18\sqrt2$}
\opt{B}{length of $AC<AT$}
\opt{C}{radius of the circumcircle is $12\sqrt2$}
\opt{D}{area of the triangle is $72\left(3+\sqrt3\right)$ sq. units.}
\end{qoptions}
\qn{9} If a focal chord of an ellipse $y^2=b^2\left(1-x^2\right)$ cuts the ellipse at the points whose eccentric angles are $\dfrac{5\pi}{12}$ and $\dfrac{23\pi}{12}$, then the value of $b$ can be : ($b<1$)
\begin{qoptions}
\opt{A}{$\dfrac{\sqrt3}2$}
\opt{B}{$\dfrac1{\sqrt3}$}
\opt{C}{$\dfrac12$}
\opt{D}{$\dfrac5{2\sqrt{10}}$}
\end{qoptions}
\qn{10} If $\log_ba=2$, $\log_cb=3$ and $a+b+c=27+4\sqrt3$ then value of :
\begin{qoptions}
\opt{A}{$\log_3(abc)=5$}
\opt{B}{$\log_3(a^2bc)=8$}
\opt{C}{$\log_3(abc^3)=6$}
\opt{D}{$\log_3(a^2bc^2)=6$}
\end{qoptions}
\sect{SECTION-III}
\qn{1} The quadratic polynomial $P(x)$ satisfy the equation $P(x)-P(x+1)=2x-1\ \forall\ x\in R$. If the greatest value of $P(x)$ in $(-\infty,0]$ is 2, then value of $P(-1)+P(1)$ is equal to
\qn{2} If number of arrangements of letters of the word ``DHARAMSHALA'' taken all at a time so that no two alike letters appear together is $\left(4^a.5^b.6^c.7^d\right)$, (where $a,b,c,d\in N$), then $a+b+c+d$ is equal to
\qn{3} Every ray of light, passing through $(1,2)$ get incident at any point $P$ lying on an elliptical surface whose eccentricity is $\dfrac{2\sqrt5}{2+\sqrt5+\sqrt{45}}$, always passes through $(3,6)$ after reflection. If $P$ is $(\sin\theta,\cos\theta)$ then value of $|3\cos\theta-6\sin\theta|$ is
\qn{4} Let $L_1$, $L_2$ and $L_3$ be the length of tangents drawn from a point $P$ to the circles $x^2+y^2=4$, $x^2+y^2-4x=0$ and $x^2+y^2-4y=0$ respectively. If $L_1^4=L_2^2L_3^2+16$, then locus of $P$ are the curves $C_1$(a straight line) and $C_2$(a circle). Let the circumcentre of the triangle formed by $C_1$ and other two lines which are inclined at an angle of $45^\circ$ with $C_1$ and is tangent to $C_2$ is $(\alpha,\beta)$, then $(\alpha+2\beta)$ equals
\qn{5} Sum of integers in range of
\par\noindent\hfill{\footnotesize$\dfrac{(4t^2+1)-\left[(4t^2-1)\cos2\theta+2t\sin2\theta\right]}{(4t^2+1)+\left[(1-4t^2)\cos2\theta+2t\sin2\theta\right]}$}\hfill\null\par
, $t\in R$, $\theta\in\left(0,\dfrac\pi2\right)$ is
\qn{6} If $a$, $b$, $c$ are the roots of the equation $x^3-2x^2+1=0$, then the value of $\left(\dfrac1a+\dfrac1{a^2}+\dfrac1b+\dfrac1{b^2}+\dfrac1c+\dfrac1{c^2}\right)$ is equal to
\qn{7} In a $\Delta ABC$, $\tan A=2$, $\tan B=\dfrac32$ and $c=2\sqrt{65}$. If circumradius of the triangle is $k$, then value of $(7k/13)$ is
\qn{8} If $2\sec^2A-\sec^4A-2\mathrm{cosec}^2A+\mathrm{cosec}^4A=\dfrac{15}4$ then $2\tan^2A$ is equal to
\end{multicols}
\end{document}
